%% =====================================================================
%%  B4-T-04-02 -- what B4 contributes to "The Favour Debt"
%%  -------------------------------------------------------------------
%%  LAYER A: APPEND-ONLY. Written once, then left alone. Material found
%%  only here sits in the `unique` slot and is protected by rule R10.
%% =====================================================================
%% @kind: source
%% @id: B4-T-04-02
%% @book: B4
%% @topic: T-04-02
%% @locator: ch. 2, pp. 43--86
%% @status: distilled
%% @unique: yes
%% @integrated: yes
%% @updated: 2026-09-19

\dossier{B4}{T-04-02}{\bookshort{B4}}{ch. 2, pp. 43--86}

\begin{distilled}
  \item The favour debt is the strongest of the compliance levers because it is enforced by visible social marking rather than by any agreement between the parties.
  \item The account's undefined balance is not a weakness of the lever but its main feature: because no amount is stated, the creditor nominates the repayment.
  \item An unsettled account is unpleasant to hold, and relieving that state is itself a motive to comply with whatever is asked.
\end{distilled}

\begin{unique}
  \item Why the debt is hard to refuse: the enforcement is public marking, not private guilt. B1 treats gratitude as an unreliable instrument; B4 explains the pressure that makes it reliable in the moment even when it is unreliable later.
  \item Discomfort as an independent driver, separate from any calculation about the transaction.
\end{unique}
